\documentclass{article}
\usepackage[margin=1in]{geometry}
\usepackage{amsmath, amssymb}

\title{Correctness Proof of the Sweep Line Algorithm for the Union Area of Rectangles}
\author{}
\date{}

\begin{document}

\maketitle

\section*{Problem Setting}

Let $R = \{R_1, R_2, \dots, R_n\}$ be a set of axis-aligned rectangles in $\mathbb{R}^2$.  
Each rectangle $R_i$ is defined by coordinates
\[
R_i = [x_i^{\min}, x_i^{\max}] \times [y_i^{\min}, y_i^{\max}]
\]
where $x_i^{\min} < x_i^{\max}$ and $y_i^{\min} < y_i^{\max}$.

The goal is to compute the area of the union
\[
A = \text{area}\left(\bigcup_{i=1}^{n} R_i\right).
\]

The sweep line algorithm processes vertical rectangle boundaries from left to right and accumulates the area contributed by vertical strips.

\section*{Definition (Event Coordinates)}

Let
\[
X = \{x_i^{\min}, x_i^{\max} \mid 1 \le i \le n\}
\]
be the set of all rectangle boundary $x$-coordinates.

Let
\[
x_1 < x_2 < \dots < x_k
\]
be the sorted sequence of distinct values in $X$.

These coordinates partition the plane into vertical strips
\[
S_j = [x_j, x_{j+1}] \quad \text{for } j = 1,2,\dots,k-1.
\]

\section*{Lemma 1 (Coverage Invariance Within a Strip)}

Within each strip $S_j$, the set of rectangles intersecting the strip is constant.

\textbf{Proof.}

Rectangles only begin or end at their vertical boundaries $x_i^{\min}$ or $x_i^{\max}$.  
Since the strip $S_j$ contains no such boundary in its interior, no rectangle starts or stops intersecting the strip within it.

Therefore, for every $x \in (x_j, x_{j+1})$, the set of rectangles intersecting the vertical line $x$ is identical.  
\hfill $\square$

\section*{Lemma 2 (Area of a Strip)}

Let $Y_j$ denote the union of the $y$-intervals of rectangles intersecting strip $S_j$.

Then the area contributed by the strip is
\[
\text{Area}(S_j) = (x_{j+1} - x_j)\cdot \text{length}(Y_j).
\]

\textbf{Proof.}

Since the set of intersecting rectangles is constant throughout the strip, the set of $y$ values covered by rectangles is the same for every $x$ in the strip.

Thus the covered region in the strip forms a Cartesian product of the strip width and the union of vertical intervals.

Therefore the area equals the strip width multiplied by the total vertical coverage length.  
\hfill $\square$

\section*{Lemma 3 (Correct Computation of Vertical Coverage)}

The interval merging procedure used in the algorithm correctly computes
\[
\text{length}\left(\bigcup_{i} [y_i^{\min}, y_i^{\max}]\right).
\]

\textbf{Proof.}

The algorithm sorts the intervals by their starting coordinate and then merges overlapping intervals.  
Whenever two intervals overlap, they are combined into a single interval whose length equals the union of the two.

Because every interval is processed exactly once and overlaps are merged, each point in the union of intervals is counted exactly once.

Therefore the algorithm computes the correct total covered length.  
\hfill $\square$

\section*{Theorem (Correctness of the Sweep Line Algorithm)}

The sweep line algorithm computes the area
\[
A = \text{area}\left(\bigcup_{i=1}^{n} R_i\right).
\]

\textbf{Proof.}

The event coordinates partition the plane into vertical strips $S_1, S_2, \dots, S_{k-1}$.

By Lemma 1, the set of rectangles intersecting each strip is constant.

By Lemma 2, the area contribution of each strip is
\[
(x_{j+1}-x_j)\cdot \text{length}(Y_j).
\]

By Lemma 3, the algorithm correctly computes the vertical coverage length $\text{length}(Y_j)$.

Since the strips are disjoint and together cover the entire region where rectangles may appear, summing the contributions of all strips gives

\[
A = \sum_{j=1}^{k-1} (x_{j+1}-x_j)\cdot \text{length}(Y_j).
\]

Thus the algorithm exactly computes the union area of all rectangles.  
\hfill $\square$

\end{document}